\pdfbookmark[1]{Of Nouns, Verbs, and Adjectives}{Of Nouns, Verbs, and Adjectives}
\subsection*{Of Nouns}

\begin{itemize}

	\item It is assumed that the reader knows about word formation and how to form plurals in Na'vi. Thus, not all plural forms are given. However, with words where the prefix causes elision plurals are given, e.g. \B{'en} (\emph{pl.} \B{men}, \B{pxen}, \B{ayen\slash en}) or where there are irregular plurals, e.g. \B{'ll\-ngo} (\emph{pl.} \B{me\-'ll\-ngo}, \B{pxe\-'ll\-ngo}, \emph{always:} \B{ay\-'ll\-ngo}) or: \B{'u} (\emph{no short plural, always:} \B{ay\-u}). Where these elisions cause pairs of homonyms, a subsequent entry will refer to the main entry, e.g. \B{pxen}\textsuperscript{\on{2}}~:~\B{'en}.

	\item This project aims to become an exhaustive reference dictionary. Thus, proper names of people, places, flora and fauna and loan words have been incorporated into the main body of the dictionary.

\end{itemize}

\subsection*{Of Verbs}

\begin{itemize}

	\item It is assumed that the reader knows about word formation and the basic concept of how Na'vi creates tense, aspect and affect in verbs by use of infixes. Thus, these infix positions in regular verbs (mono\hyp{}syllabic and di\hyp{}syllabic verbs) as well as in \B{si}\hyp{}verbs have \emph{not been marked}. However, in irregular verbs, i.e. compound verbs with elements of other word classes or two verbs, the infix positions have been indicated with the system devised by Dr.~Frommer, i.e. \B{sre\-se\-'a}, \emph{inf.}\on{23}, which gives \B{sres}‹pre-1›‹1›\B{e'}‹2›\B{a}. Additionally they are given in the familiar way from the \emph{Taronyu} Dictionary via the raised dot in the IPA (\I{[sRE.s$\cdot$E."P$\cdot$a]}).

	\item The appearance of a verb can change quite drastically through the use of infixes. To familiarize the user\slash reader with that change, the example sentences always try to cover as many different forms as possible, starting from the core verb in use, to the ‹\B{iv}› form to more complicated constructs if they could be found.
	
\end{itemize}

\subsection*{Of Adjectives}

Adjectives undergo the least changes. In their basic function they only appear in their root form given in the dictionary as the predicative adjective (\B{tute lu 'ewan}, ``The person is young'') or with an attached \B{a} as the attributive adjective (\B{tute a'ewan lu tstunwi} \emph{or} \B{'ewana tute lu tstunwi}, ``The young person is kind''). It has been attempted to find and show at least one example of each of these uses.

\pdfbookmark[1]{A Note on Spelling}{A Note on Spelling and Hyphenation}
\subsection*{A Note on Spelling and Hyphenation}

Careful thought has been given to the fact that Na'vi is a spoken language and that rules of syllabification are handled differently than in English or other languages. Prefixes, infixes and suffixes can change the appearance of a word quite drastically. A vowel can make up a whole syllable. In general, the sequence VCV (where V stands for a vowel, pseudo vowel or diphthong and C for a consonant) will be syllabified V.CV rather than VC.V. However, in instances of case endings this rule has been slightly and carefully altered to isolate the suffix as well as the root word. The reasoning being that the root word in \B{tsafnetxonìri} (``concerning that kind of night'') would be more recognizable when given as \B{tsa-fne-txon-ìri} rather than \B{tsa-fne-txo-nì-ri}, as it would be spoken. For the same reason, some words---especially those with long vowel clusters---have been prevented from a line break to leave them recognizable, e.g. \B{tsaioiyä} (``of that adornment'') as \B{tsa-ioi-yä} rather than \B{tsa-i-o-i-yä}. 

In the same vein, \B{fayhilvanit} (``those rivers'' (\emph{patientive})) is syllabified as \B{fay-hil-van-it} rather than \B{fay-hil-va-nit} to retain the root noun \B{kilvan} (river) as visible as possible.

Agentive, patientive or dative single endings have not been treated that way to avoid illegal word boundaries, e.g. \B{ayyayol} (``the birds'' (\emph{agentive})) is still syllabified as \B{ay-ya-yol} rather than *\B{ay-ya-yo-l}.

This practice was not followed, however, with verbs affected by infixes. Infixes rearrange the core syllabification of a verb. Ignoring this would lead to illegal consonant clusters. Whereas \B{terkup} (``die'') is syllabified as \B{ter-kup}, adding infixes rearranges it---in the most extreme case---to e.g. \B{täpeykìyeverkeiup} (``cause oneself to die''), syllabified as \B{tä-pey-kì-ye-ver-ke-i-up}. 

Recognizing the root verb is a matter of practice. In extreme cases the translation should give a hint as to the core verb.


